\documentclass[12pt]{article}

\usepackage[margin=1in]{geometry}
\usepackage{fancyhdr}
\usepackage{graphicx}
\usepackage{booktabs}
\usepackage{longtable}
\usepackage{array}
\usepackage{amsmath}
\usepackage{hyperref}
\usepackage{xcolor}
\usepackage{fancyhdr}
\usepackage{graphicx}
\usepackage{setspace}
\usepackage{titlesec}
\usepackage{enumitem}

%---------------------------------------------------------------------------------
% Section Formatting (optional: smaller spacing before sections)
%---------------------------------------------------------------------------------
\titlespacing\section{0pt}{*2}{*1}
\titlespacing\subsection{0pt}{*1.5}{*0.8}
\titlespacing\subsubsection{0pt}{*1.2}{*0.5}

%---------------------------------------------------------------------------------
% Cover Page
%---------------------------------------------------------------------------------
\pagestyle{fancy}
\fancyhf{}
\fancyhead[L]{USC Facilities: Rain Barrel Water Pump System}
\fancyfoot[C]{\thepage}

\begin{document}
\onehalfspacing

\title{
    \vspace{2in}
    \textbf{EMCH428-002} \\[1ex]
    \textbf{Test Plan} \\[1ex]
    \textbf{Sponsor: USC Facilities} \\[1ex]
    \vspace{2in}
}

\author{
    \textbf{Team Members:} \\[2ex]
    Tanner Harmon\\[1ex]
    Presston McConnell\\[1ex]
    William Wenzel\\[1ex]
    Jacob Vaughn\\[2ex]
}

\date{}

\maketitle
\thispagestyle{fancy}

\newpage

\setcounter{page}{2} % Start page numbering from the second page

%----------------------------------------------------------------------------------------
% 1. PRODUCT MISSION STATEMENT
%----------------------------------------------------------------------------------------
\section*{Product Mission Statement}
\textbf{Mission:} 
Provide the university with a self-contained, solar-compatible pump system that pressurizes water from a 1000-gallon rain tank, delivering at least 12~gal/min at usable pressure for irrigation while reducing manual labor and emphasizing sustainability.

\noindent
This system should:
\begin{itemize}
    \item Supply consistent water flow and pressure for garden irrigation.
    \item Operate reliably using existing solar/battery infrastructure, minimizing reliance on grid power.
    \item Decrease volunteer/manual watering tasks while promoting water conservation.
\end{itemize}


%----------------------------------------------------------------------------------------
% 2. NEEDS / ENGINEERING CHARACTERISTICS
%----------------------------------------------------------------------------------------
\section*{List of Needs and Engineering Characteristics (E.C.)}

\vspace{1em}
\begin{tabular}{p{0.5cm}p{5.5cm}p{8cm}}
\toprule
\textbf{\#} & \textbf{Need / E.C.} & \textbf{Brief Description} \\
\midrule
1 & Flow Rate & Must deliver at least 12 gal/min (target) for efficient watering. \\
2 & Pressure & Must maintain 40--60 psi at the hose outlet for typical garden sprinklers. \\
3 & Power Source & Integrate with existing solar/battery or minimal AC usage. \\
4 & Durability & System must withstand outdoor elements (UV, rain, partial freezing). \\
5 & Ease of Assembly & Must be easy to assemble by student volunteers with minimal tools. \\
6 & Cost & Must remain within budget constraints for parts; low annual operating cost. \\
7 & Safety & Must have safe wiring (GFCI) and over-pressure protection. \\
\bottomrule
\end{tabular}


%----------------------------------------------------------------------------------------
% 3. METRICS MATRIX (HOUSE OF QUALITY)
%----------------------------------------------------------------------------------------
\section*{Up-to-date Metrics Matrix / House of Quality (HOQ)}

\vspace{1em}
\begin{center}
\begin{tabular}{>{\centering\arraybackslash}p{2.8cm} | >{\centering\arraybackslash}p{1.5cm}
                                      >{\centering\arraybackslash}p{1.5cm}
                                      >{\centering\arraybackslash}p{1.5cm}
                                      >{\centering\arraybackslash}p{1.5cm}
                                      >{\centering\arraybackslash}p{1.5cm}
                                      >{\centering\arraybackslash}p{1.5cm}}
\toprule
\textbf{Needs} $\downarrow$ 
   & \rotatebox{90}{\parbox{2.0cm}{\centering Flow Rate}} 
   & \rotatebox{90}{\parbox{2.0cm}{\centering Pressure}} 
   & \rotatebox{90}{\parbox{2.0cm}{\centering Power Source}} 
   & \rotatebox{90}{\parbox{2.0cm}{\centering Durability}} 
   & \rotatebox{90}{\parbox{2.0cm}{\centering Ease of Assembly}} 
   & \rotatebox{90}{\parbox{2.0cm}{\centering Cost}} \\ 
\midrule
N1: Deliver water quickly    
   & 9 & 9 & 3 & 1 & 1 & 3 \\ 
N2: Handle sprinklers 
   & 3 & 9 & 1 & 3 & 1 & 1 \\ 
N3: Limit manual labor   
   & 9 & 1 & 9 & 3 & 3 & 3 \\ 
N4: Survive weather  
   & 1 & 3 & 1 & 9 & 1 & 3 \\
N5: Easy to build 
   & 3 & 3 & 3 & 1 & 9 & 3 \\
N6: Stay on budget 
   & 3 & 3 & 3 & 3 & 3 & 9 \\
\bottomrule
\end{tabular}
\end{center}

\noindent
\textbf{Legend:}  
- “9” = Strong Relationship;  
- “3” = Moderate Relationship;  
- “1” = Weak Relationship.  

This matrix maps how each engineering characteristic helps fulfill each stated customer need.


%----------------------------------------------------------------------------------------
% 4. NUMBERED LIST OF SPECIFICATIONS
%----------------------------------------------------------------------------------------
\section*{Numbered Specifications (Target \& Fallback Values)}
Below is our up-to-date list of measurable specifications mapped to the E.C. above.

\begin{enumerate}
    \item \textbf{Flow Rate:} 
    \begin{itemize}
        \item Target: 12--16 gal/min
        \item Fallback: 8--12 gal/min
    \end{itemize}

    \item \textbf{Pressure:}
    \begin{itemize}
        \item Target: 40--60 psi
        \item Fallback: 20--40 psi
    \end{itemize}

    \item \textbf{Power Source (Energy Efficiency):}
    \begin{itemize}
        \item Target: $<$\,\$50 per year for electricity
        \item Fallback: $<$\,\$150 per year
    \end{itemize}

    \item \textbf{Durability (Outdoor Survival):}
    \begin{itemize}
        \item Target: 10 years lifespan with minimal maintenance
        \item Fallback: 5+ years lifespan
    \end{itemize}

    \item \textbf{Ease of Assembly:}
    \begin{itemize}
        \item Target: \textless{}5 hours total build time
        \item Fallback: \textless{}8 hours total build time
    \end{itemize}

    \item \textbf{Cost (Capital \& Operating):}
    \begin{itemize}
        \item Target: \textless{}\$2,000 initial build, \textless{}\$100/year maintenance
        \item Fallback: \textless{}\$2,500 initial
    \end{itemize}

    \item \textbf{Safety (Electrical/Pressure):}
    \begin{itemize}
        \item Target: GFCI, relief valve, safe wiring, & no injuries
        \item Fallback: Standard breaker, manual check
    \end{itemize}
\end{enumerate}


%----------------------------------------------------------------------------------------
% 5. TEST PLANS FOR EACH NEED/SPEC
%----------------------------------------------------------------------------------------
\newpage\section*{Test Plans for Each Assessed Specification}

Below, we provide the \textbf{NUMBER} (from the specs above) and the \textbf{Parameter Name}, followed by the required test plan details.

%---------------------------------------
% 5.1 Flow Rate
%---------------------------------------
\subsection*{Spec \#1: Flow Rate}
\textbf{Parameter Name:} Flow Rate (gal/min)

\noindent
\textbf{Materials Required:}
\begin{itemize}
    \item 5-gallon bucket (or similar known volume container)
    \item Stopwatch
    \item The installed pump system with a standard 3/4'' garden hose
\end{itemize}

\noindent
\textbf{Procedure:}
\begin{enumerate}
    \item Ensure the rain barrel has sufficient water volume (at least 50 gallons) to allow continuous pumping.
    \item Connect the system’s outlet hose to the 5-gallon bucket (or place the hose end inside it).
    \item Turn on the pump; simultaneously start the stopwatch.
    \item Time how long it takes to fill the 5-gallon bucket.
    \item Repeat for at least 3 trials to improve accuracy.
    \item If feasible, test additional flow rates with a sprinkler attached (to replicate actual usage).
\end{enumerate}

\noindent
\textbf{Data to Collect:}
\begin{itemize}
    \item Time in seconds to fill 5 gallons (Trial 1, Trial 2, Trial 3, etc.)
    \item Observations about flow consistency or cavitation
\end{itemize}

\noindent
\textbf{Data Analysis:}
\begin{itemize}
    \item Flow Rate $= \dfrac{\text{Volume (gal)}}{\text{Time (min)}}$
    \item Compare average flow rate vs. target range (12--16 gal/min).
    \item If the average < 8 gal/min, the fallback spec is also not met.
\end{itemize}

\vspace{1em}
%---------------------------------------
% 5.2 Pressure
%---------------------------------------
\subsection*{Spec \#2: Pressure}
\textbf{Parameter Name:} Outlet Pressure (psi)

\noindent
\textbf{Materials Required:}
\begin{itemize}
    \item Inline pressure gauge (0--100 psi range)
    \item Standard 3/4'' garden hose
    \item The installed pump system
\end{itemize}

\noindent
\textbf{Procedure:}
\begin{enumerate}
    \item Attach the inline pressure gauge between the pump outlet and the garden hose.
    \item Start the pump under normal operating conditions (with the main valve open).
    \item Record the stabilized pressure reading on the gauge while water is flowing.
    \item Observe if the pressure switch cycles on/off and note any fluctuations.
    \item Repeat at least 2 more times, especially if the system has an automated cutoff or tank.
\end{enumerate}

\noindent
\textbf{Data to Collect:}
\begin{itemize}
    \item Steady-state outlet pressure readings (psi)
    \item On/off cycle pressures (e.g., cut-in/cut-out)
\end{itemize}

\noindent
\textbf{Data Analysis:}
\begin{itemize}
    \item Compare measured average pressure to target 40--60 psi range.
    \item If average is below 20 psi or above 60 psi, note potential failures or adjustments needed.
\end{itemize}

\vspace{1em}
%---------------------------------------
% 5.3 Power Source (Energy Efficiency)
%---------------------------------------
\newpage\subsection*{Spec \#3: Power/Energy Efficiency}
\textbf{Parameter Name:} Annual Electricity Cost (\$/year)

\noindent
\textbf{Materials Required:}
\begin{itemize}
    \item Clamp-on ammeter or watt-meter to measure pump’s electrical draw
    \item Timer or usage log
    \item Local electricity rate (\$/kWh)
\end{itemize}

\noindent
\textbf{Procedure:}
\begin{enumerate}
    \item Measure the pump's power draw (watts or amps \& voltage) under typical operating conditions.
    \item Estimate daily runtime based on irrigation schedule (e.g., 30 minutes/day).
    \item Use your local electricity rate (e.g., \$0.12/kWh) to approximate daily cost.
    \item Scale daily cost to annual cost.
\end{enumerate}

\noindent
\textbf{Data to Collect:}
\begin{itemize}
    \item Pump current draw (A), voltage (V), or direct wattage from watt-meter.
    \item Pump runtime (hours/day).
    \item Electricity rate \$/kWh.
\end{itemize}

\noindent
\textbf{Data Analysis:}
\begin{itemize}
    \item Annual Cost $\approx \text{(Watts)} \times \text{(hours/day)} \times 365 \times \frac{\text{\$}}{1000\,\text{W}\cdot\text{hr}}$
    \item Compare result to target $<$ \$50/year (or fallback $<$ \$150/year).
\end{itemize}

\vspace{1em}
%---------------------------------------
% 5.4 Durability (Outdoor Survival)
%---------------------------------------
\newpage\subsection*{Spec \#4: Durability (Outdoor Survival)}
\textbf{Parameter Name:} 10-year Lifespan / Weather Resistance

\noindent
\textbf{Materials Required:}
\begin{itemize}
    \item Visual inspection checklist for structural integrity (housing, piping)
    \item Simulated weather exposure (e.g., water spray test, if relevant)
    \item Basic mechanical tools to check for loosened fasteners
\end{itemize}

\noindent
\textbf{Procedure:}
\begin{enumerate}
    \item After system is assembled, perform a once-a-semester or once-a-year inspection:
        \begin{itemize}
            \item Check for leaks, cracks, or rust around pump, tank, and piping.
            \item Inspect housing for water infiltration or rot.
        \end{itemize}
    \item Optionally, do a short-term accelerated test by spraying water or leaving the system in harsh sunlight for a day, then re-check for problems.
\end{enumerate}

\noindent
\textbf{Data to Collect:}
\begin{itemize}
    \item Observations of any structural damage, leaks, or corrosion
    \item Condition of protective paint or sealants on wood or metal components
\end{itemize}

\noindent
\textbf{Data Analysis:}
\begin{itemize}
    \item Rate the condition from 1--5 or pass/fail on the checklist each year.
    \item If major failure occurs before year 5, fallback spec is not met.
\end{itemize}

\vspace{1em}

%---------------------------------------
% 5.7 Safety (Electrical/Pressure)
%---------------------------------------
\newpage\subsection*{Spec \#7: Safety (Electrical/Pressure)}
\textbf{Parameter Name:} GFCI / Relief Valve Verification

\noindent
\textbf{Materials Required:}
\begin{itemize}
    \item GFCI outlet tester
    \item Visual inspection of relief valve setting
    \item Pressure gauge to confirm valve opens at set point (optional)
\end{itemize}

\noindent
\textbf{Procedure:}
\begin{enumerate}
    \item Plug the pump system into a GFCI outlet. Use a GFCI tester to ensure it trips properly.
    \item Verify the presence of a relief valve or pressure regulator on the tank or piping.
    \item If possible, raise system pressure to test the relief valve opening at the correct psi (e.g., 70 psi).
\end{enumerate}

\noindent
\textbf{Data to Collect:}
\begin{itemize}
    \item GFCI test result (pass/fail)
    \item Relief valve rated psi vs. actual open psi
\end{itemize}

\noindent
\textbf{Data Analysis:}
\begin{itemize}
    \item The system must pass the GFCI test to meet the safety target.
    \item Relief valve must open below 75 psi to prevent over-pressure conditions.
\end{itemize}


%----------------------------------------------------------------------------------------
% END
%----------------------------------------------------------------------------------------
\end{document}
